% Source-aligned English for printed p.128.
\SourcePage{128}

\LayerTarget{tohoku:section:1.4:p4}
\noindent the canonical surjection $B\to B/B'$. Dually, one
defines the direct image $u(A')$ of a quotient $A'$ of $A$; it is a quotient
of $B$. All the usual formal properties are proved for these notions.

\LayerTarget{tohoku:section:1.4:exact-sequence}
Finally, recall that a pair $A\mathrel{\mathop{\to}^{u}}B \mathrel{\mathop{\to}^{v}}C$ of two consecutive morphisms is called \textit{exact} if
$\mathrm{Ker}\,v=\mathrm{Im}\,u$, whence, more generally, the notion of an
\textit{exact sequence} of morphisms. For a sequence $0\to A\to B\to C$ to be exact, it is necessary and sufficient that, for every
$X\in\mathbf{C}$, the following sequence of homomorphisms of abelian groups be exact:

\[
  0\to\mathrm{Hom}(X,A)\to\mathrm{Hom}(X,B)
  \to\mathrm{Hom}(X,C).
\]

\noindent There is a dual criterion for a sequence
$C\to B\to A\to0$ to be exact. For a
sequence $0\to A'\to A\to A''\to0$ to be exact, it is necessary and sufficient that $u$ be a
monomorphism and that $v$ be a generalized cokernel of it.

\LayerTarget{tohoku:section:1.4:exact-functor}
Let $F$ be a covariant functor from an abelian category $\mathbf{C}$ to another such
category $\mathbf{C}'$. Following the terminology introduced in [6], we say that $F$ is
a \textit{half-exact functor} (respectively, \textit{left exact}, respectively,
\textit{right exact}) if, for every exact sequence $0\to A'\to A\to A''\to0$ in $\mathbf{C}$, the corresponding sequence of morphisms $0\to F(A')\to F(A)\to F(A'') \to0$ is exact at $F(A)$ (respectively, at $F(A)$ and $F(A')$,
respectively, at $F(A)$ and $F(A'')$). $F$ is called an \textit{exact functor}
if it is both left and right exact, i.e. if it carries an exact sequence of the
preceding type to an exact sequence; it then carries \textit{every} exact
sequence to an exact sequence. If $F$ is left exact, it carries an exact
sequence $0\to A\to B\to C$ to an exact sequence $0\to F(A)\to F(B)\to F(C)$; there is a dual statement for right-exact functors. If $F$ is a
contravariant functor, it is said to be half-exact (respectively, left exact,
etc.) if it has that property as a covariant functor from $\mathbf{C}^\circ$ to $\mathbf{C}'$.
The composite of covariant left-exact (respectively, right-exact) functors is
of the same type. We refer to [6] for further sorites of this kind and for the
study of exactness properties of multifunctors. As an important example, note
that $\mathrm{Hom}(A,B)$ is an additive bifunctor on $\mathbf{C}\times\mathbf{C}$, with values
in the abelian category of abelian groups, contravariant in $A$ and covariant
in $B$, and left exact with respect to both arguments (that is, a left-exact
bifunctor in the terminology of [6]).

\LayerTarget{tohoku:section:1.5}
\noindent\textbf{1.5. Infinite sums and products.}
In certain constructions we shall need the existence and certain properties
of infinite direct sums and direct products. Here, in increasing order of
strength, are the most frequently used axioms.

\LayerTarget{tohoku:section:1.5:AB3}
\noindent AB 3) \textit{For every family $(A_i)_{i\in I}$ of objects
of $\mathbf{C}$, the direct sum of the $A_i$ (cf. No.~1) exists.}

\LayerTarget{tohoku:section:1.5:p1}
This axiom implies that, for every family of subobjects $A_i$ of an object
$A\in\mathbf{C}$, the supremum of the $A_i$ exists: it suffices to take the image of
the direct sum $\oplus A_i$ under the morphism whose components are
the canonical injections $A_i\to A$. We have seen that the direct
sum of an arbitrary family of surjective morphisms is surjective (No.~1); in
fact, one even sees that the functor $(A_i)_{i\in I}\to\displaystyle\bigoplus_{i\in I}A_i$, defined on the ``product category'' $\mathbf{C}^I$, and taking values

\TohokuContinuation{129}
